% Part: first-order-logic
% Chapter: models-theories
% Section: size-of-structures

\documentclass[../../../include/open-logic-section]{subfiles}

\begin{document}

\olfileid{fol}{mat}{siz}

\olsection{\printtoken{P}{structure} పరిమాణాన్ని వ్యక్తం చేయడం}

\begin{explain}
ఒక భాషలోని తార్కికేతర సంకేతాలను ఉపయోగించకుండానే వ్యక్తం చేయగల కొన్ని
నిర్మాణ ధర్మాలు ఉన్నాయి. ఉదాహరణకు, \tetoken{నిర్మాణం}{!!{structure}} యొక్క \tetoken{వ్యక్తి క్షేత్రంలో}{!!{domain}}
కనీసం, గరిష్ఠంగా, లేదా సరిగ్గా~$n$ \tetoken{మూలకాలు}{!!{element}s}
ఉన్నప్పుడు మాత్రమే ఆ \tetoken{ఒక నిర్మాణంలో}{!!a{structure}} సత్యమయ్యే \tetoken{వాక్యాలు}{!!{sentence}s} ఉన్నాయి.
\end{explain}

\begin{prop}
\tetoken{వాక్యం}{!!{sentence}}
\begin{multline*}
  !A_{\ge n} \ident \lexists[x_1][\lexists[x_2][\dots\lexists[x_n][{}]]]\\
  \begin{aligned}
  (\eq/[x_1][x_2] \land {}
  \eq/[x_1][x_3] \land \eq/[x_1][x_4] \land \dots \land \eq/[x_1][x_n] \land {}\\
\eq/[x_2][x_3] \land \eq/[x_2][x_4] \land \dots \land {} \eq/[x_2][x_n] \land {} \\
\vdots\\
\eq/[x_{n-1}][x_n])
\end{aligned}
\end{multline*}
\tetoken{ఒక నిర్మాణం}{!!a{structure}}~$\Struct M$లో సత్యం కావడం, $\Domain M$లో కనీసం
$n$ \tetoken{మూలకాలు}{!!{element}s} ఉండటానికి అవసరమైన మరియు సరిపడే షరతు. తత్ఫలితంగా,
$\Domain M$లో గరిష్ఠంగా~$n$ \tetoken{మూలకాలు}{!!{element}s} ఉన్నప్పుడు మాత్రమే
$\Sat{M}{\lnot !A_{\ge n+1}}$.

\end{prop}

\begin{prop}
\tetoken{వాక్యం}{!!{sentence}}
\begin{multline*}
!A_{= n} \ident \lexists[x_1][\lexists[x_2][\dots\lexists[x_n][{}]]] \\
\begin{aligned}
  (\eq/[x_1][x_2] \land {}
  \eq/[x_1][x_3] \land \eq/[x_1][x_4] \land \dots \land \eq/[x_1][x_n] \land {}\\
\eq/[x_2][x_3] \land \eq/[x_2][x_4] \land \dots \land {} \eq/[x_2][x_n] \land {} \\
\vdots\\
\eq/[x_{n-1}][x_n] \land {} \\
\lforall[y][(\eq[y][x_1] \lor \dots \lor \eq[y][x_n]]))
\end{aligned}
\end{multline*}
\tetoken{ఒక నిర్మాణం}{!!a{structure}}~$\Struct M$లో సత్యం కావడం, $\Domain M$లో సరిగ్గా
$n$ \tetoken{మూలకాలు}{!!{element}s} ఉండటానికి అవసరమైన మరియు సరిపడే షరతు.
\end{prop}

\begin{prop}
ఒక \tetoken{నిర్మాణం}{!!{structure}} కింది సమితికి నమూనా అయితే, అప్పుడు మాత్రమే అది అనంతమైనది:
\[
\{!A_{\ge 1}, !A_{\ge 2}, !A_{\ge 3}, \dots \}.
\]
\end{prop}

$\Domain M$ అనంతమైనప్పుడు మాత్రమే~$\Struct M$లో సత్యమయ్యే ఒకే శుద్ధ
తార్కిక వాక్యం లేదు. అయితే అనంత నమూనాలు మాత్రమే గల, తార్కికేతర
\tetoken{విధేయ సంకేతాలతో}{!!{predicate}s} కూడిన \tetoken{వాక్యాలను}{!!{sentence}s} ఇవ్వవచ్చు (అయినా ప్రతి అనంత
\tetoken{నిర్మాణం}{!!{structure}} వాటికి నమూనా కాదు). పరిమిత నిర్మాణంగా ఉండే ధర్మాన్ని,
\tetoken{లెక్కించలేని}{!!{nonenumerable}} నిర్మాణంగా ఉండే ధర్మాన్ని అనంత \tetoken{వాక్యాలు}{!!{sentence}s}
సమితితో కూడా వ్యక్తం చేయలేం. ఈ వాస్తవాలు సంహతత్వ, లొవెన్‌హైమ్--స్కోలెమ్
సిద్ధాంతాల నుంచి అనుగమిస్తాయి.

\end{document}
